En este capítulo se expone el \textbf{contexto} general del trabajo, los conceptos tecnológicos indispensables que motivan a su desarrollo y, finalmente, la estructura en la que se organiza el resto del documento.

\section{Contexto y motivación}
La preservación de la memoria parlamentaria y la transparencia en la gestión pública por parte de los representantes se encuentran principalmente en los Diarios de Sesiones del Parlamento de Canarias, convirtiéndolos en una valiosa fuente documental. Durante décadas, estas actas han registrado las intervenciones, debates y acuerdos legislativos del archipiélago. Actualmente, el acceso y la consulta a este gran conjunto de ficheros son facilitados por el \gls{iatext} \cite{iatext_web} a través de herramientas de recuperación especializadas como el buscador textual \textbf{\gls{disecan}}.

No obstante, la evolución tecnológica y las crecientes necesidades de investigadores y ciudadanos demandan sistemas de recuperación de información más intuitivos y precisos disponibles para el público. Los buscadores tradicionales, basados puramente en coincidencias \textit{booleanas} o expresiones regulares (\textit{Regex}), presentan ciertas limitaciones al no ser capaces de comprender ni el contexto semántico ni la intención detrás del mensaje que presenta el usuario. Esta limitación es el principal motivo para el desarrollo de este trabajo, cuyo propósito es investigar e implementar arquitecturas modernas basadas en \gls{iagen} para conseguir una versión del sistema de búsqueda más contextual y eficaz.

\section{Conceptos básicos: Hacia el RAG Híbrido}
Para superar los inconvenientes de los buscadores tradicionales, este proyecto se basa en la arquitectura de Generación Aumentada por Recuperación (\gls{rag}) \cite{lewis2020rag}. Esta técnica permite a los \gls{llm} generar respuestas precisas consultando bases de datos específicas y cerradas como fuentes de verdad, mitigando así el riesgo de generar información falsa o \textit{alucinaciones} \cite{ji2023hallucination}.


Concretamente, se propone el diseño de un sistema \textbf{\gls{rag} Híbrido}. Esta solución integra la precisión de las bases de datos relacionales (\gls{sql}), ideales para filtrar por metadatos institucionales (como el nombre del orador, la fecha o la legislatura), con la flexibilidad que ofrecen las bases de datos vectoriales, capaces de realizar búsquedas por similitud semántica. Todo ello siendo ejecutado en una infraestructura local para garantizar la \textbf{soberanía, seguridad y privacidad} de los datos utilizados.

Como resultado, este trabajo aporta un prototipo funcional de asistente conversacional para el \gls{disecan}, compuesto por un motor de búsqueda híbrido y una interfaz web interactiva diseñada para la comodidad y usabilidad por parte del usuario final.

\section{Organización del documento}
El resto del documento está integrado por los siguientes capítulos, estructurados conforme a las directrices de la Escuela de Ingeniería Informática:

\begin{itemize}
    \item El \textbf{Capítulo 2} expone el estado actual del problema, los antecedentes tecnológicos y los objetivos específicos planteados inicialmente para el desarrollo del trabajo.
    \item El \textbf{Capítulo 3} detalla las aportaciones del proyecto, la justificación de las competencias adquiridas durante el grado y el alineamiento con los \gls{ods}.
    \item El \textbf{Capítulo 4} describe la metodología aplicada (\gls{crisp-dm}) y las fases técnicas de diseño, desarrollo e implementación de la arquitectura híbrida.
    \item El \textbf{Capítulo 5} presenta los resultados obtenidos tras la experimentación, validación y despliegue del sistema propuesto.
    \item El \textbf{Capítulo 6} recoge las conclusiones derivadas del proyecto, evaluando la consecución de objetivos, y propone posibles líneas de trabajo futuro.
    \item Finalmente, el \textbf{Capítulo 7} incluye consideraciones adicionales como el marco normativo y la legislación aplicable a los sistemas de Inteligencia Artificial en el sector público.
\end{itemize}